\cleardoublepage

\newrefsection

\chapter{外文翻译}

\sectionnonum{摘要}

以下为论文《面向实时辐射场渲染的三维高斯泼溅》的译文\cite{kerbl2023}。辐射场方法最近彻底改变了对由多张照片或视频采集得到的场景进行新视角合成的方式。然而，要获得高视觉质量，仍然需要训练和渲染代价高昂的神经网络，而近来的更快方法又不可避免地在速度和质量之间进行权衡。对于无界且完整的场景，而不是孤立物体，以及在 \(1080p\) 分辨率下的渲染，目前还没有任何方法能够达到实时显示速率。我们引入了三个关键要素，使我们能够在保持具有竞争力的训练时间的同时达到最先进的视觉质量，并且更重要的是，在 \(1080p\) 分辨率下实现高质量实时（\(\ge 30\) 帧每秒）的新视角合成。第一，从相机标定过程中产生的稀疏点出发，我们使用三维高斯来表示场景。这种表示在场景优化时保留了连续体辐射场的理想性质，同时避免了空区域中的不必要计算。第二，我们对三维高斯执行交替的优化与密度控制，尤其优化各向异性协方差，以获得对场景的精确表示。第三，我们提出了一种快速、考虑可见性的渲染算法，它支持各向异性泼溅，并同时加速训练和实现实时渲染。我们在多个经典数据集上展示了最先进的视觉质量和实时渲染能力。

\section{引言}

网格和点是最常见的三维场景表示，因为它们是显式的，并且非常适合基于 GPU/CUDA 的快速光栅化。与此相对，近来的神经辐射场方法建立在连续场景表示之上，通常通过体积光线步进优化多层感知机来对采集场景进行新视角合成。类似地，到目前为止最高效的辐射场方案也建立在连续表示之上，它们通过插值存储在体素网格、哈希网格或点中的数值来工作。尽管这些方法的连续性有利于优化，但渲染时所需的随机采样代价高昂，并且会带来噪声。我们提出了一种结合两者优点的新方法：我们的三维高斯表示可以在保持最先进视觉质量和具有竞争力训练时间的同时进行优化，而基于分块的泼溅方案则保证了在若干公开数据集上以 \(1080p\) 分辨率实现具备最先进质量的实时渲染。

我们的目标，是针对由多张照片采集得到的场景实现实时渲染，并且在典型真实场景上，使这种表示的优化时间与先前最高效的方法一样快。近来的方法已经实现了较快训练，但很难达到当前最先进 NeRF 方法 Mip-NeRF360 所取得的视觉质量，而后者最多需要 \(48\) 小时训练。那些速度快但质量较低的辐射场方法，依赖场景情况可以实现交互式渲染速度，即每秒 \(10\) 到 \(15\) 帧，但在高分辨率下仍然达不到实时渲染。

我们的方案建立在三个主要组成部分之上。首先，我们引入三维高斯作为一种灵活而富有表达力的场景表示。我们从与先前类 NeRF 方法相同的输入开始，即由运动恢复结构法校准得到的相机，并使用作为该过程副产物的稀疏点云来初始化三维高斯集合。与多数需要多视图立体数据的点基础方案不同，我们仅使用 SfM 点作为输入，就能得到高质量结果。需要说明的是，对于 NeRF Synthetic 数据集，即使采用随机初始化，我们的方法也能达到高质量。我们表明，三维高斯是一个非常好的选择，因为它是一种可微的体表示，同时又可以通过投影到二维并进行标准透明度混合而被非常高效地光栅化，其图像形成模型与 NeRF 等价。我们方法的第二个组成部分，是对三维高斯属性的优化，也就是三维位置、透明度 \(\alpha\)、各向异性协方差以及球谐系数，并在优化过程中穿插执行自适应密度控制步骤，在其中我们会添加、偶尔也会移除三维高斯。这个优化过程产生了一种相当紧凑、非结构化且精确的场景表示。我们方法的第三个也是最后一个组成部分，是一种实时渲染方案，它使用快速 GPU 排序算法，并受到了基于分块光栅化工作的启发。得益于三维高斯表示，我们能够执行遵循可见性顺序的各向异性泼溅，并通过跟踪所需数量的已排序泼溅，完成快速而准确的反向传播。

总之，本文的贡献如下：
\begin{itemize}
    \item 引入各向异性的三维高斯，作为一种高质量、非结构化的辐射场表示。
    \item 提出一种三维高斯属性的优化方法，并与自适应密度控制交替进行，从而为采集场景构建高质量表示。
    \item 提出一种面向 GPU 的快速可微渲染方法。该方法考虑可见性，允许各向异性泼溅，并支持快速反向传播，从而实现高质量新视角合成。
\end{itemize}

我们在既有公开数据集上的结果表明，我们能够从多视图采集中优化三维高斯，并取得与先前最优隐式辐射场方法相当甚至更好的质量。我们同样能够获得与最快方法相近的训练速度和质量，并且更重要的是，首次在高质量新视角合成中实现了实时渲染。

\section{相关工作}

\subsection{传统场景重建与渲染}

我们首先简要回顾传统重建方法，然后讨论点基础渲染和辐射场工作，并说明它们之间的相似性。辐射场是一个很大的研究领域，因此这里只关注与本文直接相关的工作。

最早的新视角合成方法建立在光场之上，最初是稠密采样，随后发展到允许非结构化采集。运动恢复结构法的出现开启了一个全新的方向，使得一组照片可以用来合成新视角。SfM 在相机标定期间会估计出一个稀疏点云，这在最初主要用于三维空间的简单可视化。随后，多视图立体逐渐发展出令人印象深刻的完整三维重建算法，并推动了若干视图合成算法的发展。所有这些方法都会将输入图像重投影并混合到新的视角相机中，并使用几何信息来引导这种重投影。这些方法在许多情况下都能取得优秀结果，但通常无法完全从未重建区域恢复，也难以处理 MVS 产生不存在几何时的过度重建。近来的神经渲染算法大幅减轻了这类伪影，并且避免了将所有输入图像存放在 GPU 上所带来的巨大开销，因此在大多数方面都优于这些传统方法。

\subsection{神经渲染与辐射场}

深度学习技术很早就被用于新视角合成；卷积神经网络曾被用于估计混合权重，或者用于纹理空间方案。大多数这类方法的主要缺点在于它们依赖 MVS 几何；此外，使用卷积神经网络进行最终渲染经常会导致时间上的闪烁。

用于新视角合成的体表示由 Soft3D 开始，随后出现了将深度学习技术与体积光线步进结合的方法，这些方法建立在连续可微的密度场之上来表示几何。由于需要大量样本来查询体数据，体积光线步进渲染的代价非常高。NeRF 引入了重要性采样和位置编码来提高质量，但使用了较大的多层感知机，从而对速度产生了负面影响。NeRF 的成功引发了大量后续工作，这些工作主要通过引入正则化策略来处理质量和速度问题；当前在新视角合成图像质量方面的最先进方法是 Mip-NeRF360。尽管其渲染质量极高，但训练和渲染时间仍然非常长；而我们在提供快速训练和实时渲染的同时，能够达到甚至在某些情况下超过这一质量。

最近的方法主要通过三种设计选择来关注更快的训练和渲染：使用空间数据结构存储随后在体积光线步进中进行插值的特征，使用不同的编码方式，以及调整多层感知机的容量。这些方法包括多种空间离散化变体、码本以及哈希表等编码方式，使得可以使用更小的多层感知机，甚至完全不使用神经网络。其中最著名的方法包括 InstantNGP，它使用哈希网格和占据网格加速计算，并使用较小的多层感知机来表示密度和外观；以及 Plenoxels，它使用稀疏体素网格来插值连续密度场，并且能够完全去掉神经网络。二者都依赖球谐函数：前者直接用它表示方向相关效应，后者用它对颜色网络的输入进行编码。尽管这些方法给出了非常出色的结果，但它们仍然可能难以有效表示空区域，这在一定程度上取决于场景和采集方式。此外，图像质量在很大程度上仍受用于加速的结构化网格限制，而渲染速度则受限于体积光线步进中需要查询的大量样本。我们使用的非结构化、显式、对 GPU 友好的三维高斯，在不引入神经组件的情况下实现了更快的渲染速度和更好的质量。

\subsection{点基础渲染与辐射场}

点基础方法能够高效渲染不连通、非结构化的几何采样，也就是点云。在最简单的形式下，点采样渲染以固定大小对一个非结构化点集进行光栅化，它可以利用图形接口原生支持的点类型，也可以利用 GPU 上的并行软件光栅化。点采样渲染虽然忠实于原始数据，但会产生孔洞，引起混叠，并且是严格不连续的。高质量点基础渲染的经典工作，通过对范围大于一个像素的点元进行泼溅，例如圆盘、椭圆体或曲面元，来解决这些问题。

近来，人们对可微点基础渲染技术表现出浓厚兴趣。点被增强了神经特征，并使用卷积神经网络进行渲染，从而实现了快速甚至实时的视图合成；然而，这些方法仍然依赖 MVS 作为初始几何，因此也继承了它的伪影，最明显的是在无纹理区域、强反光区域或细薄结构等困难场景中的过重建或欠重建。

点基础透明度混合与 NeRF 风格的体渲染本质上共享同一种图像形成模型。具体来说，颜色 \(C\) 由沿光线的体渲染给出：
\[
C=\sum_{i=1}^{N} T_i\left(1-\exp\left(-\sigma_i\delta_i\right)\right)c_i,
\qquad
T_i=\exp\left(-\sum_{j=1}^{i-1}\sigma_j\delta_j\right).
\]
其中，沿光线以间隔 \(\delta_i\) 采样密度 \(\sigma_i\)、透射率 \(T_i\) 和颜色 \(c_i\)。这一表达式可以改写为
\[
C=\sum_{i=1}^{N} T_i\alpha_i c_i,
\qquad
\alpha_i=1-\exp(-\sigma_i\delta_i),\quad
T_i=\prod_{j=1}^{i-1}(1-\alpha_j).
\]
而典型的神经点基础方法通过混合覆盖该像素的有序点来计算像素颜色：
\[
C=\sum_{i=1}^{N} c_i\alpha_i\prod_{j=1}^{i-1}(1-\alpha_j),
\]
其中 \(c_i\) 是每个点的颜色，而 \(\alpha_i\) 则由协方差为 \(\Sigma\) 的二维高斯与学习得到的逐点透明度相乘得到。

从上述公式可以清楚地看到，二者的图像形成模型是相同的。然而，它们的渲染算法却非常不同。NeRF 是一种连续表示，它以隐式方式表示空区域和占据区域；为了在公式中找到样本，需要进行代价高昂的随机采样，这会带来噪声和计算开销。与之相反，点是一种非结构化、离散的表示，它足够灵活，允许像 NeRF 一样对几何进行创建、销毁和位移。已有工作已经表明，可以通过优化透明度和位置来做到这一点，同时避免完整体表示的缺陷。

Pulsar 实现了快速的球体光栅化，这启发了我们基于分块和排序的渲染器。不过，根据上面的分析，我们希望保留对已排序泼溅进行传统透明度混合的近似，以拥有体表示的优点：我们的光栅化遵循可见性顺序，而他们的方法与顺序无关。此外，我们对一个像素中的所有泼溅都反向传播梯度，并且支持各向异性泼溅。这些因素共同促成了结果中的高视觉质量。除此之外，前面提到的方法还使用卷积神经网络进行渲染，这会带来时间不稳定性。尽管如此，Pulsar 和 ADOP 的渲染速度仍然促使我们发展出自己的快速渲染方案。

还有一种较新的方法使用点和径向基函数来表示辐射场。它们在优化过程中采用点裁剪和加密技术，但仍然使用体积光线步进，因此无法达到实时显示速率。在人体表演捕获领域，三维高斯曾被用于表示采集到的人体；近来，它们也被与体积光线步进结合，用于视觉任务。神经体元方法也在相近背景下被提出。这些工作启发了我们选择三维高斯作为场景表示，但它们关注的是重建和渲染单个孤立物体，因此场景的深度复杂度较低。与此相对，我们对各向异性协方差的优化、优化与密度控制的交替进行以及渲染时高效的深度排序，使我们能够处理包含背景的完整复杂场景，包括室内和室外场景，以及具有较大深度复杂度的情况。

\section{总体概览}

我们方法的输入是一组静态场景图像，以及由 SfM 校准得到的对应相机参数；这一过程还会附带生成一个稀疏点云。我们从这些点出发，创建一组三维高斯。每个高斯由位置，也就是均值，协方差矩阵以及透明度 \(\alpha\) 定义，这使得我们可以采用一种非常灵活的优化机制。由此得到的三维场景表示相当紧凑，这在一定程度上是因为高度各向异性的体泼溅可以紧凑地表示精细结构。辐射场的方向相关外观分量，也就是颜色，则按照标准做法使用球谐函数来表示。随后，我们的算法通过一系列对三维高斯参数的优化步骤来构建辐射场表示，这些参数包括位置、协方差、\(\alpha\) 和球谐系数，并在其中穿插执行用于自适应控制高斯密度的操作。整体流程如图~\ref{fig:3dgs-overview} 所示。

我们方法效率的关键，是基于分块的光栅化器。它通过快速排序来实现各向异性泼溅的透明度混合，并遵守可见性顺序。我们的快速光栅化器还通过跟踪累积的 \(\alpha\) 值来实现快速反向传播，而且不限制能够接收梯度的高斯数量。

\begin{figure}[H]
    \centering
    \includegraphics[width=0.95\textwidth]{translation/3dgs-fig2-overview.png}
    \caption{方法总览。优化从稀疏 SfM 点云开始，并创建一组三维高斯；随后对高斯参数进行优化，并自适应地控制其密度。}
    \label{fig:3dgs-overview}
\end{figure}

\section{可微三维高斯泼溅}

我们的目标，是从一组没有法向的稀疏 SfM 点出发，优化得到一种能够实现高质量新视角合成的场景表示。要做到这一点，我们需要一种图元：它既继承可微体表示的性质，同时又必须是非结构化和显式的，以便能够实现非常快的渲染。我们选择三维高斯，因为它是可微的，并且可以很容易地投影为二维泼溅，从而支持快速的透明度混合渲染。

我们的表示与先前使用二维点的方法有相似之处。那些方法假设每个点都是带有法向的小平面圆盘。考虑到 SfM 点极其稀疏，很难估计法向；同样地，去优化由这类估计得到的、噪声很大的法向也会非常困难。因此，我们将几何建模为一组三维高斯，而不需要法向。我们的高斯由定义在世界坐标系中的完整三维协方差矩阵 \(\Sigma\) 给出，其中心位于点，也就是均值 \(\mu\) 处：
\[
G(\mathbf{x})=\exp\!\left(-\frac{1}{2}(\mathbf{x}-\mu)^{T}\Sigma^{-1}(\mathbf{x}-\mu)\right).
\]
在我们的混合过程中，这个高斯还会乘以 \(\alpha\)。

然而，为了进行渲染，我们需要把三维高斯投影到二维。已有工作展示了如何将这种投影变换到图像空间。给定观察变换 \(W\)，相机坐标系中的协方差矩阵 \(\Sigma'\) 可表示为
\[
\Sigma' = J W \Sigma W^{T} J^{T},
\]
其中 \(J\) 是投影变换仿射近似的雅可比矩阵。相关工作还表明，如果忽略 \(\Sigma'\) 的第三行和第三列，我们就能得到一个 \(2\times2\) 的方差矩阵，它的结构和性质与从带法向的平面点出发时得到的结果相同。

一种显然的做法，是直接优化协方差矩阵 \(\Sigma\)，从而得到表示辐射场的三维高斯。然而，只有当协方差矩阵是半正定时，它才具有物理意义。我们对全部参数采用梯度下降进行优化，而这种优化并不容易被约束为始终产生有效矩阵；更新步骤和梯度也很容易生成无效的协方差矩阵。

因此，我们选择了一种更直观、但表达能力等价的优化表示。三维高斯的协方差矩阵 \(\Sigma\) 可以看作对一个椭球体形状的描述。给定缩放矩阵 \(S\) 和旋转矩阵 \(R\)，我们可以得到相应的协方差矩阵：
\[
\Sigma = R S S^{T} R^{T}.
\]
为了使这两个因素能够独立优化，我们分别存储它们：用三维向量 \(s\) 表示缩放，用四元数 \(q\) 表示旋转。它们可以很容易地转换成各自对应的矩阵并组合起来，同时对 \(q\) 进行归一化以得到有效的单位四元数。

为了避免训练过程中自动微分带来的显著开销，我们显式推导了全部参数的梯度。这样一种适合优化的各向异性协方差表示，使我们能够优化三维高斯，使之适应采集场景中不同形状的几何，从而得到一种相当紧凑的表示。优化后的高斯形状示意如图~\ref{fig:3dgs-anisotropic-shapes} 所示。

\begin{figure}[H]
    \centering
    \includegraphics[width=0.72\textwidth]{translation/3dgs-fig3-shapes.png}
    \caption{优化后的三维高斯形状示意。将高斯收缩后可以更清楚地看到其各向异性形状，这种表示能够较为紧凑地贴合复杂几何。}
    \label{fig:3dgs-anisotropic-shapes}
\end{figure}

\section{带自适应密度控制的优化}

\subsection{优化}

这种方法的核心，是优化步骤。该步骤会构建出一组稠密的三维高斯，以精确表示场景，从而用于自由视点合成。除了位置 \(\mu\)、透明度 \(\alpha\) 和协方差 \(\Sigma\) 之外，我们还优化表示每个高斯颜色 \(c\) 的球谐系数，以便正确捕获场景的视角相关外观。这些参数的优化，与控制高斯密度以更好表示场景的步骤交替进行。

优化建立在连续的渲染与比较迭代之上。我们在采集数据集的训练视图上渲染当前结果，并将生成图像与真实图像进行比较。由于三维到二维投影的歧义，几何不可避免地可能被错误放置。因此，我们的优化既需要能够创建几何，也需要在几何被错误放置时销毁或移动它。三维高斯协方差参数的质量，对于表示的紧凑性至关重要，因为较大的均匀区域可以由少量较大的各向异性高斯来表示。

我们使用随机梯度下降技术进行优化，并充分利用标准 GPU 加速框架以及为某些操作添加自定义 CUDA 内核的能力。特别是，我们的快速光栅化过程是优化效率的关键，因为它是优化中最主要的计算瓶颈。我们对 \(\alpha\) 使用 Sigmoid 激活函数，将其约束在 \([0,1)\) 范围内并获得平滑梯度；对协方差尺度使用指数激活函数，原因也类似。我们将初始协方差矩阵估计为一个各向同性高斯，其各轴长度等于与最近三个点距离均值。我们采用与 Plenoxels 类似的标准指数衰减调度策略，但仅用于位置参数。损失函数为 \(L_1\) 与 D-SSIM 的组合：
\[
L = (1-\lambda)L_1 + \lambda L_{D\text{-}SSIM},
\]
在我们的所有测试中都取 \(\lambda=0.2\)。

\subsection{高斯的自适应控制}

我们从 SfM 产生的初始稀疏点集出发，然后应用本文的方法，对高斯的数量以及单位体积上的密度进行自适应控制，从而使我们能够从一个稀疏的初始高斯集合，过渡到一个更稠密、能够更好表示场景且参数正确的高斯集合。在优化预热之后，我们每 \(100\) 次迭代执行一次加密，并移除任何基本透明的高斯，也就是透明度 \(\alpha\) 小于阈值 \(\tau_\alpha\) 的高斯。这里的高斯密度当然不应与 NeRF 文献中的体密度 \(\sigma\) 混淆。

我们对高斯的自适应控制需要去填充空白区域。它既关注存在几何特征缺失的区域，也关注高斯在场景中覆盖过大范围的区域，而后者通常对应于过度重建。我们观察到，这两种情况都会产生较大的视空间位置梯度。直观地看，这是因为这些区域尚未被很好地重建，而优化过程正试图移动这些高斯来纠正它。

由于这两种情况都适合执行加密，我们会对那些视空间位置梯度平均模长高于阈值 \(\tau_{\text{pos}}\) 的高斯执行加密。在我们的测试中，该阈值设为 \(0.0002\)。这两种加密情形的示意如图~\ref{fig:3dgs-density} 所示。接下来我们给出该过程的细节。

对于位于欠重建区域中的小高斯，我们需要去覆盖必须被创建出来的新几何。对此，更合适的做法是克隆这些高斯，即直接创建一个相同大小的拷贝，并沿着位置梯度的方向移动它。

另一方面，对于处在高方差区域中的大高斯，则需要把它们分裂为更小的高斯。我们用两个新的高斯替换这样的高斯，并将它们的尺度除以 \(\phi=1.6\)，这个数值是我们通过实验确定的。我们还利用原始三维高斯作为概率密度函数来采样，从而初始化它们的位置。

在第一种情况下，我们识别并处理的是既需要增加系统总体覆盖体积、又需要增加高斯数量的情形；而在第二种情况下，我们保持总体体积不变，但增加高斯数量。与其他体表示类似，我们的优化也可能在输入相机附近陷入漂浮物问题；在我们的情形下，这可能会导致高斯密度不合理地增加。为了有效抑制高斯数量的增长，我们每隔 \(N=3000\) 次迭代就把 \(\alpha\) 的值设得接近于零。随后，优化会在真正需要的地方重新增大这些高斯的 \(\alpha\)，同时允许我们的裁剪策略移除那些 \(\alpha\) 小于阈值 \(\tau_\alpha\) 的高斯。高斯可能会收缩或增大，并与其他高斯发生显著重叠，但我们会周期性地移除那些在世界空间中过大、以及在视空间中占据很大范围的高斯。该策略使我们能够较好地控制高斯总数。与其他方法不同，我们模型中的高斯始终是欧氏空间中的图元；我们不需要对远处或巨大的高斯进行空间压缩、扭曲或投影处理。

\begin{figure}[H]
    \centering
    \includegraphics[width=0.75\textwidth]{translation/3dgs-fig4-density.png}
    \caption{我们的自适应高斯加密方案。上排对应欠重建情形：当小尺度几何覆盖不足时，我们克隆相应高斯。下排对应过度重建情形：当小尺度几何被一个大泼溅表示时，我们将其一分为二。}
    \label{fig:3dgs-density}
\end{figure}

\section{面向高斯的快速可微光栅化}

我们的目标，是同时获得快速的整体渲染和快速排序，以便允许近似的透明度混合，包括对各向异性泼溅的透明度混合，并避免先前工作中存在的、可接收梯度泼溅数量的硬限制。

为实现这一目标，我们设计了一种面向高斯泼溅的基于分块的光栅化器，其灵感来自近来的软件光栅化方法。该方法会先对整幅图像中的图元进行预排序，从而避免逐像素排序的代价，而逐像素排序正是先前透明度混合方案的瓶颈。我们的快速光栅化器允许对任意数量的已混合高斯进行高效反向传播，同时只引入很低的额外内存消耗，也就是每个像素仅需常数级额外开销。考虑到前面完成的二维投影，这一光栅化流程是完全可微的，并且能够像先前的二维泼溅方法一样，对各向异性泼溅进行光栅化。

我们的方法首先将屏幕划分为 \(16\times16\) 的小块，然后继续对三维高斯执行视锥裁剪以及逐块裁剪。具体而言，我们只保留其 \(99\%\) 置信区间与视锥相交的高斯。此外，我们使用一个保护带，直接剔除那些位置极端的高斯，也就是均值靠近近裁剪面且远在视锥之外的高斯，因为为这些高斯计算投影后的二维协方差会不稳定。随后，我们按照每个高斯所覆盖的块数实例化高斯，并为每个实例分配一个同时编码视空间深度和块编号的键值。然后，我们使用一次快速 GPU 基数排序，按照这些键对高斯进行排序。需要注意的是，这里并不存在额外的逐像素点排序，混合完全依据这一次初始排序进行。因此，在某些配置下，我们的透明度混合是近似的。不过，当泼溅逐渐接近单个像素大小时，这些近似会变得可以忽略不计。我们发现，这一选择在不产生可见伪影的前提下，显著提升了训练和渲染性能。

在高斯排序之后，我们通过确定某个给定块上按照深度排序后的首个和末个条目，为每个块生成一个列表。进行光栅化时，我们为每个块启动一个线程块。每个线程块首先协同地把成批高斯加载到共享内存中，然后对给定像素以前向到后向的顺序遍历列表，累积颜色和 \(\alpha\) 值，从而同时最大化数据加载、共享和处理过程中的并行性。当某个像素中的 \(\alpha\) 达到目标饱和度后，对应线程便停止。在固定间隔上会检查块中的各个线程；一旦所有像素都达到饱和，也就是 \(\alpha\) 趋近于 \(1\)，整个块的处理就终止。

在光栅化过程中，\(\alpha\) 的饱和是唯一的停止准则。与先前工作不同，我们不限制能够接收梯度更新的混合图元数量。我们坚持这一性质，是为了使方法能够处理具有任意且变化深度复杂度的场景，并准确学习这些场景，而不必借助针对场景的超参数调节。因此，在反向传播时，我们必须恢复前向传播中每个像素所混合点的完整序列。一种做法，是在全局内存中为每个像素存储任意长度的混合点列表。为了避免由此带来的动态内存管理开销，我们选择再次遍历逐块列表；我们可以复用前向过程中已经得到的已排序高斯数组和块范围。为了便于梯度计算，我们现在按从后到前的顺序遍历它们。

这种遍历从影响该块中任一像素的最后一个点开始，点再次以协作方式被加载到共享内存中。此外，某个像素只有在点的深度小于或等于它在前向过程中最后一个对其颜色有贡献的点的深度时，才会开始执行代价较高的重叠测试和点处理。前面提到的梯度计算，需要知道原始混合过程中每一步累积的不透明度值。与其在反向过程中显式遍历一个不断减小的不透明度列表，我们只需存储前向过程结束时的总累积不透明度，就可以恢复这些中间值。具体来说，每个点在前向过程中会存储最终累积的不透明度 \(A\)；在从后到前的遍历中，我们将这个值依次除以各点的 \(\alpha\)，从而得到梯度计算所需的系数。

\section{实现、结果与评估}

\subsection{实现}

我们使用 PyTorch 框架在 Python 中实现了该方法，并为光栅化编写了自定义 CUDA 内核。快速基数排序使用 NVIDIA CUB 排序例程。我们还构建了一个基于 SIBR 的交互式查看器，并使用它来测量实际获得的帧率。

为了保持稳定性，我们在较低分辨率下对计算进行预热。具体来说，我们从缩小四倍分辨率的图像开始优化，并在第 \(250\) 次和第 \(500\) 次迭代之后各上采样一次。球谐系数的优化对角度信息缺失非常敏感。对于典型的类 NeRF 采集方式，也就是围绕中心物体的整个半球都有照片观测，优化可以很好工作。然而，如果采集中存在缺失的角区域，例如采集场景角落或执行从内向外的采集，那么球谐零阶分量，也就是基础色或漫反射颜色，可能会被优化出完全错误的值。为了解决这一问题，我们先只优化零阶分量，然后每隔 \(1000\) 次迭代引入一个新的球谐带，直到全部 \(4\) 个带都被表示出来。

\subsection{结果与评估}

我们在总计 \(13\) 个真实场景以及 Blender 合成数据集上测试了算法。具体而言，我们在 Mip-NeRF360 的全部场景、Tanks\&Temples 数据集中的两个场景，以及 Hedman 等人提供的两个场景上测试了方法。我们选择的场景具有非常不同的采集方式，同时涵盖了受限的室内场景和大型无界室外环境。在全部实验中，我们都使用同一组超参数配置。所有结果都在 A6000 GPU 上报告。

在真实场景上，就质量而言，当前最先进的方法是 Mip-NeRF360，因此我们将其作为质量基准。同时，我们也与两个最新的快速 NeRF 方法，也就是 InstantNGP 和 Plenoxels 进行比较。对于数据集划分，我们采用与 Mip-NeRF360 建议相同的训练集与测试集划分方式，即每隔 \(8\) 张照片取一张作为测试图像，以便在误差指标上进行一致而有意义的比较。我们使用文献中最常见的 PSNR、LPIPS 和 SSIM 作为评价指标。除 Mip-NeRF360 在其自有数据集上的结果外，表中的数值均来自我们对作者代码的自行运行；对于该部分，我们直接采用原论文中的数字，以避免对当前最先进结果产生混淆。我们同时报告平均训练时间、渲染速度以及用于存储优化后参数的显存占用。对于 InstantNGP，我们报告一个训练 \(35\mathrm{K}\) 次迭代的基础配置，以及作者建议的稍大网络配置；对于我们的方法，则报告 \(7\mathrm{K}\) 和 \(30\mathrm{K}\) 两种配置。

\begin{table}[H]
    \centering
    \caption{我们的方法与先前工作的定量比较，结果在三个数据集上计算。带匕首符号的结果直接取自原论文，其余结果均来自作者的实验。}
    \scriptsize
    \resizebox{\textwidth}{!}{
    \begin{tabular}{l|ccc|ccc|ccc}
        \toprule
        \multirow{2}{*}{方法} & \multicolumn{3}{c|}{Mip-NeRF360} & \multicolumn{3}{c|}{Tanks\&Temples} & \multicolumn{3}{c}{Deep Blending} \\
        & PSNR$\uparrow$ & SSIM$\uparrow$ & LPIPS$\downarrow$ & PSNR$\uparrow$ & SSIM$\uparrow$ & LPIPS$\downarrow$ & PSNR$\uparrow$ & SSIM$\uparrow$ & LPIPS$\downarrow$ \\
        \midrule
        Plenoxels & 23.08 & 0.626 & 0.463 & 21.08 & 0.719 & 0.379 & 23.06 & 0.795 & 0.510 \\
        INGP-Base & 25.30 & 0.671 & 0.371 & 21.72 & 0.723 & 0.330 & 23.62 & 0.797 & 0.423 \\
        INGP-Big & 25.59 & 0.699 & 0.331 & 21.92 & 0.745 & 0.305 & 24.96 & 0.817 & 0.390 \\
        Mip-NeRF360 & 27.69 & 0.792 & 0.237 & 22.22 & 0.759 & 0.257 & 29.40 & 0.901 & 0.245 \\
        Ours-7K & 25.60 & 0.770 & 0.279 & 21.20 & 0.767 & 0.280 & 27.78 & 0.875 & 0.317 \\
        Ours-30K & 27.21 & 0.815 & 0.214 & 23.14 & 0.841 & 0.183 & 29.41 & 0.903 & 0.243 \\
        \bottomrule
    \end{tabular}}
\end{table}

在许多情况下，\(7\mathrm{K}\) 次迭代时的质量已经相当不错。训练时间在不同数据集之间有所变化，因此我们分别报告。表中的结果表明，我们的完全收敛模型在质量上与 Mip-NeRF360 不相上下，某些情况下甚至略好一些；需要注意的是，在相同硬件上，后者的平均训练时间为 \(48\) 小时，而我们仅需 \(35\) 到 \(45\) 分钟，并且其渲染时间约为每帧 \(10\) 秒。经过 \(5\) 到 \(10\) 分钟训练后，我们就能达到与 InstantNGP 和 Plenoxels 相近的质量；而继续训练则使我们能够达到最先进质量，这是其他快速方法做不到的。对于 Tanks\&Temples，我们在与基础版 InstantNGP 近似的训练时间下取得了相近的质量。

\begin{figure}[H]
    \centering
    \includegraphics[width=\textwidth,trim=70 315 70 80,clip]{translation/3dgs-page-07.png}
    \caption{我们的方法、先前方法以及测试视图真实图像之间的比较。图中从上到下分别为 Bicycle、Garden、Stump、Counter、Room、Playroom、DrJohnson、Truck 和 Train。箭头和局部放大框标出了不易直接察觉的质量差异。}
    \label{fig:3dgs-qualitative}
\end{figure}

上图展示了与先前渲染方法的可视化比较，所示为一张未参与训练的测试视图。图中的我们的方法采用 \(30\mathrm{K}\) 次迭代的结果。可以看到，在某些情况下，即使是 Mip-NeRF360 也仍然存在我们能够避免的伪影，例如 Bicycle 和 Stump 中植被区域的模糊，以及 Room 中墙面的模糊。我们的办法倾向于在良好覆盖区域中，即使从远处观察，也能保留视觉细节，而先前方法并不总是如此。

除了真实场景之外，我们还在 Blender 合成数据集上评估了方法。这些场景提供了完备的视图集合、有限的场景尺度以及精确的相机参数。在这类场景中，即使采用随机初始化，我们也能够达到最先进结果：我们从包围场景边界体积内的 \(100\mathrm{K}\) 个均匀随机高斯开始训练，随后快速、自动地将其裁剪为约 \(6\) 到 \(10\mathrm{K}\) 个有意义的高斯。在 \(30\mathrm{K}\) 次迭代后，训练完成的模型最终达到每个场景约 \(200\) 到 \(500\mathrm{K}\) 个高斯。下表给出了在白色背景设置下与先前方法的 PSNR 比较。训练完成后的合成场景渲染速度可达到每秒 \(180\) 到 \(300\) 帧。

\begin{table}[H]
    \centering
    \caption{Synthetic NeRF 数据集上的 PSNR 分数。我们从 \(100\mathrm{K}\) 个随机初始化点开始，比较方法的数值取自各自原论文。}
    \label{tab:3dgs-syn-psnr}
    \scriptsize
    \resizebox{\textwidth}{!}{
    \begin{tabular}{lccccccccc}
        \toprule
        方法 & Mic & Chair & Ship & Materials & Lego & Drums & Ficus & Hotdog & Avg. \\
        \midrule
        Plenoxels & 33.26 & 33.98 & 29.62 & 29.14 & 34.10 & 25.35 & 31.83 & 36.81 & 31.76 \\
        INGP-Base & 36.22 & 35.00 & 31.10 & 29.78 & 36.39 & 26.02 & 33.51 & 37.40 & 33.18 \\
        Mip-NeRF & 36.51 & 35.14 & 30.41 & 30.71 & 35.70 & 25.48 & 33.29 & 37.48 & 33.09 \\
        Point-NeRF & 35.95 & 35.40 & 30.97 & 29.61 & 35.04 & 26.06 & 36.13 & 37.30 & 33.30 \\
        Ours-30K & 35.36 & 35.83 & 30.80 & 30.00 & 35.78 & 26.15 & 34.87 & 37.72 & 33.32 \\
        \bottomrule
    \end{tabular}}
\end{table}

与先前的显式场景表示相比，我们优化中使用的各向异性高斯能够以更少的参数来建模复杂形状。为说明这一点，作者将本文方法与高度紧凑的点基础模型进行了比较。从对方通过空间雕刻和前景掩码得到的初始点云出发，经过 \(2\) 到 \(4\) 分钟优化后，我们就能够达到甚至超过其报告的 PSNR，而且仅使用了约四分之一的点数，平均模型大小为 \(3.8\) MB，而对方为 \(9\) MB。

\subsection{消融实验}

我们把不同贡献和算法选择逐一分离，并构造了一组实验来测量它们的效果。具体来说，我们测试以下几个方面：使用 SfM 进行初始化、我们的加密策略、各向异性协方差、允许无限数量的泼溅接收梯度，以及球谐函数的使用。每一种选择的定量效果见下表。

\begin{table}[H]
    \centering
    \caption{消融实验的 PSNR 分数。该实验中，作者将各场景的高分辨率输入图像手动下采样到与其他实验一致的渲染分辨率。}
    \scriptsize
    \resizebox{\textwidth}{!}{
    \begin{tabular}{lcccccccc}
        \toprule
        方法 & Truck-5K & Garden-5K & Bicycle-5K & Truck-30K & Garden-30K & Bicycle-30K & Average-5K & Average-30K \\
        \midrule
        Limited-BW & 14.66 & 22.07 & 20.77 & 13.84 & 22.88 & 20.87 & 19.16 & 19.19 \\
        RandomInit & 16.75 & 20.90 & 19.86 & 18.02 & 22.19 & 21.05 & 19.17 & 20.42 \\
        No-Split & 18.31 & 23.98 & 22.21 & 20.59 & 26.11 & 25.02 & 21.50 & 23.90 \\
        No-SH & 22.36 & 25.22 & 22.88 & 24.39 & 26.59 & 25.08 & 23.48 & 25.35 \\
        No-Clone & 22.29 & 25.61 & 22.15 & 24.82 & 27.47 & 25.46 & 23.35 & 25.91 \\
        Isotropic & 22.40 & 25.49 & 22.81 & 23.89 & 27.00 & 24.81 & 23.56 & 25.23 \\
        Full & 22.71 & 25.82 & 23.18 & 24.81 & 27.70 & 25.65 & 23.90 & 26.05 \\
        \bottomrule
    \end{tabular}}
\end{table}

首先，我们评估从 SfM 点云初始化三维高斯的重要性。在这一消融中，作者在一个大小为输入相机包围盒三倍的立方体内进行均匀采样。结果表明，即使没有 SfM 点，我们的方法也不会完全失败，而是主要在背景部分退化。在训练视图覆盖不足的区域，随机初始化会产生更多无法通过优化去除的漂浮伪影。另一方面，Synthetic NeRF 数据集没有这种行为，因为它没有背景，并且受到输入相机的良好约束。

接下来，作者评估了两种加密方法，也就是第 5 节中描述的克隆和分裂策略。作者分别禁用其中一种，而保持方法其余部分不变。结果表明，分裂大高斯对于背景的良好重建非常重要，而克隆小高斯而不是分裂它们，则能够带来更好、更快的收敛，尤其是在场景中出现细长结构时。

作者还评估了当只允许最前面的 \(N\) 个点接收梯度，并跳过其后的梯度计算时，是否能够像 Pulsar 所暗示的那样，在不牺牲质量的前提下换取速度。在该实验中，作者取 \(N=10\)，这是 Pulsar 默认值的两倍，但这种严重近似会导致优化不稳定。对于 Truck 场景，质量在 PSNR 上下降了约 \(11\) dB。

另一个重要的算法选择，是对三维高斯完整协方差矩阵的优化。为说明这一点，作者进行了一组消融实验，在其中移除了各向异性，只优化一个控制三轴共同半径的标量值。结果表明，各向异性显著改善了三维高斯与表面对齐的能力，从而在保持相同点数的同时获得更高的渲染质量。

最后，球谐函数的使用提高了整体 PSNR 分数，因为它们能够补偿视角相关效应。

\subsection{局限性}

我们的方法并非没有局限。在场景没有被良好观测的区域，我们会出现伪影；在这些区域中，其他方法也同样表现困难。尽管各向异性高斯具有前面所述的许多优点，我们的方法仍然可能生成拉长的伪影或斑驳状高斯；先前的方法在这些情况下也会遇到困难。

当优化产生较大的高斯时，我们还会偶尔出现跳变伪影；这通常发生在具有视角相关外观的区域。造成这些跳变伪影的一个原因，是光栅化器中通过保护带对高斯进行简单剔除。若使用更有原则的裁剪方法，这些伪影会有所缓解。另一个因素，是我们当前使用的简单可见性算法，它可能导致高斯突然切换深度次序和混合次序。未来可以通过抗锯齿来处理这一问题。此外，我们目前并未对优化施加任何正则化；若加入正则化，也将有助于缓解未观测区域和跳变伪影。

尽管在完整评估中我们使用了相同的超参数，但早期实验表明，在非常大的场景中，可能需要减小位置学习率才能收敛。虽然与先前的点基础方法相比，我们已经非常紧凑，但内存消耗仍然显著高于基于 NeRF 的方案。在大型场景训练期间，我们未经优化的原型实现其峰值 GPU 显存占用可能超过 \(20\) GB。不过，这一数值可以通过更加仔细的底层实现大幅降低。渲染训练完成的场景需要足够的 GPU 显存来存储完整模型，而光栅化器还会额外占用约 \(30\) 到 \(500\) MB，具体取决于场景规模和图像分辨率。

\section{讨论与结论}

我们提出了第一种真正能够在多种场景和采集方式下，实现高质量辐射场实时渲染的方法，同时训练时间又与先前最快的方法具有竞争力。

我们选择三维高斯作为图元，使其在优化过程中保留了体渲染的性质，同时又能够直接支持基于泼溅的快速光栅化。我们的工作表明，与广泛接受的看法相反，连续表示并不是实现快速且高质量辐射场训练的必要条件。

我们的大部分训练时间，大约 \(80\%\)，花在 Python 代码中，因为我们在 PyTorch 中实现了方法，以便其他人更容易使用。只有光栅化步骤被实现为优化过的 CUDA 内核。我们预计，如果像 InstantNGP 那样把其余优化过程也全部迁移到 CUDA 中，那么在强调性能的应用里还可以进一步显著加速。

我们还展示了建立在实时渲染原则之上的重要性，也就是充分利用 GPU 的能力以及软件光栅化流水线架构的速度。这些设计选择，是训练和实时渲染都能获得高性能的关键，并使我们在性能上相较于先前的体积光线步进方法具有竞争优势。

未来如果能够利用这些高斯执行采集场景的网格重建，将会非常有意义。除了网格广泛使用所带来的实际价值之外，这也将帮助我们更好理解我们的方法在体表示和表面表示连续谱中的确切位置。

总之，我们提出了第一个面向辐射场的实时渲染方案。它的渲染质量能够匹配先前那些代价高昂但效果最好的方法，而训练时间则与现有最快方案处于同一量级。

\section{外文参考文献}

为保证文献著录信息的准确性，以下按原论文中的英文著录信息录入，并保持原文顺序。

\begingroup
\small
\sloppy
\begin{enumerate}[leftmargin=2em, itemsep=0.3em, topsep=0.4em, parsep=0pt]
    \item Kara-Ali Aliev, Artem Sevastopolsky, Maria Kolos, Dmitry Ulyanov, and Victor Lempitsky. 2020. Neural Point-Based Graphics. In \textit{Computer Vision -- ECCV 2020: 16th European Conference, Glasgow, UK, August 23--28, 2020, Proceedings, Part XXII}. 696--712.

    \item Jonathan T. Barron, Ben Mildenhall, Matthew Tancik, Peter Hedman, Ricardo Martin-Brualla, and Pratul P. Srinivasan. 2021. Mip-NeRF: A multiscale representation for anti-aliasing neural radiance fields. In \textit{Proceedings of the IEEE/CVF International Conference on Computer Vision}. 5855--5864.

    \item Jonathan T. Barron, Ben Mildenhall, Dor Verbin, Pratul P. Srinivasan, and Peter Hedman. 2022. Mip-NeRF 360: Unbounded Anti-Aliased Neural Radiance Fields. In \textit{Proceedings of the IEEE/CVF Conference on Computer Vision and Pattern Recognition}.

    \item Sebastien Bonopera, Jerome Esnault, Siddhant Prakash, Simon Rodriguez, Theo Thonat, Mehdi Benadel, Gaurav Chaurasia, Julien Philip, and George Drettakis. 2020. SIBR: A System for Image Based Rendering. \url{https://gitlab.inria.fr/sibr/sibr_core}.

    \item Mario Botsch, Alexander Hornung, Matthias Zwicker, and Leif Kobbelt. 2005. High-Quality Surface Splatting on Today's GPUs. In \textit{Proceedings of the Second Eurographics / IEEE VGTC Conference on Point-Based Graphics (SPBG '05)}. Eurographics Association, Goslar, DEU, 17--24.

    \item Chris Buehler, Michael Bosse, Leonard McMillan, Steven Gortler, and Michael Cohen. 2001. Unstructured lumigraph rendering. In \textit{Proceedings of SIGGRAPH}.

    \item Gaurav Chaurasia, Sylvain Duchene, Olga Sorkine-Hornung, and George Drettakis. 2013. Depth synthesis and local warps for plausible image-based navigation. \textit{ACM Transactions on Graphics (TOG)} 32, 3 (2013), 1--12.

    \item Anpei Chen, Zexiang Xu, Andreas Geiger, Jingyi Yu, and Hao Su. 2022b. TensoRF: Tensorial Radiance Fields. In \textit{European Conference on Computer Vision (ECCV)}.

    \item Zhiqin Chen, Thomas Funkhouser, Peter Hedman, and Andrea Tagliasacchi. 2022a. MobileNeRF: Exploiting the Polygon Rasterization Pipeline for Efficient Neural Field Rendering on Mobile Architectures. \textit{arXiv preprint arXiv:2208.00277} (2022).

    \item Ricardo L. de Queiroz and Philip A. Chou. 2016. Compression of 3D point clouds using a region-adaptive hierarchical transform. \textit{IEEE Transactions on Image Processing} 25, 8 (2016), 3947--3956.

    \item Martin Eisemann, Bert De Decker, Marcus Magnor, Philippe Bekaert, Edilson de Aguiar, Naveed Ahmed, Christian Theobalt, and Anita Sellent. 2008. Floating textures. In \textit{Computer Graphics Forum}, Vol. 27. Wiley Online Library, 409--418.

    \item John Flynn, Ivan Neulander, James Philbin, and Noah Snavely. 2016. DeepStereo: Learning to predict new views from the world's imagery. In \textit{Proceedings of the IEEE Conference on Computer Vision and Pattern Recognition}.

    \item Alex Fridovich-Keil, Matthew Tancik, Qinhong Chen, Benjamin Recht, and Angjoo Kanazawa. 2022. Plenoxels: Radiance Fields without Neural Networks. In \textit{Proceedings of the IEEE/CVF Conference on Computer Vision and Pattern Recognition}.

    \item Stephan J. Garbin, Marek Kowalski, Matthew Johnson, Jamie Shotton, and Julien Valentin. 2021. FastNeRF: High-Fidelity Neural Rendering at 200FPS. In \textit{Proceedings of the IEEE/CVF International Conference on Computer Vision (ICCV)}. 14346--14355.

    \item Michael Goesele, Noah Snavely, Brian Curless, Hugues Hoppe, and Steven M. Seitz. 2007. Multi-view stereo for community photo collections. In \textit{Proceedings of the IEEE International Conference on Computer Vision}.

    \item Steven J. Gortler, Radek Grzeszczuk, Richard Szeliski, and Michael F. Cohen. 1996. The lumigraph. In \textit{Proceedings of the 23rd Annual Conference on Computer Graphics and Interactive Techniques}. 43--54.

    \item Markus Gross and Hanspeter Pfister (Eds.). 2011. \textit{Point-Based Graphics}. Elsevier.

    \item Jeff P. Grossman and William J. Dally. 1998. Point Sample Rendering. In \textit{Rendering Techniques}.

    \item Peter Hedman, Julien Philip, True Price, Jan-Michael Frahm, George Drettakis, and Gabriel Brostow. 2018. Deep blending for free-viewpoint image-based rendering. \textit{ACM Transactions on Graphics (TOG)} 37, 6 (2018).

    \item Peter Hedman, Tobias Ritschel, George Drettakis, and Gabriel Brostow. 2016. Scalable Inside-Out Image-Based Rendering. \textit{ACM Transactions on Graphics (SIGGRAPH Asia Conference Proceedings)} 35, 6 (December 2016). \url{http://www-sop.inria.fr/reves/Basilic/2016/HRDB16}.

    \item Peter Hedman, Pratul P. Srinivasan, Ben Mildenhall, Jonathan T. Barron, and Paul Debevec. 2021. Baking Neural Radiance Fields for Real-Time View Synthesis. In \textit{Proceedings of the IEEE/CVF International Conference on Computer Vision}.

    \item Philipp Henzler, Niloy J. Mitra, and Tobias Ritschel. 2019. Escaping Plato's cave: 3D shape from adversarial rendering. In \textit{Proceedings of the IEEE/CVF International Conference on Computer Vision}. 9984--9993.

    \item Arno Knapitsch, Jaesik Park, Qian-Yi Zhou, and Vladlen Koltun. 2017. Tanks and Temples: Benchmarking large-scale scene reconstruction. \textit{ACM Transactions on Graphics (TOG)} 36, 4 (2017), 1--13.

    \item Georgios Kopanas, Thomas Leimk{\"u}hler, Gilles Rainer, Cl{\'e}ment Jambon, and George Drettakis. 2022. Neural Point Catacaustics for Novel-View Synthesis of Reflections. \textit{ACM Transactions on Graphics (SIGGRAPH Asia Conference Proceedings)} 41, 6 (2022), 201. \url{http://www-sop.inria.fr/reves/Basilic/2022/KLRJD22}.

    \item Georgios Kopanas, Julien Philip, Thomas Leimk{\"u}hler, and George Drettakis. 2021. Point-Based Neural Rendering with Per-View Optimization. \textit{Computer Graphics Forum} 40, 4 (2021), 29--43. \url{https://doi.org/10.1111/cgf.14339}.

    \item Samuli Laine and Tero Karras. 2011. High-performance software rasterization on GPUs. In \textit{Proceedings of the ACM SIGGRAPH Symposium on High Performance Graphics}. 79--88.

    \item Christoph Lassner and Michael Zollh{\"o}fer. 2021. Pulsar: Efficient Sphere-Based Neural Rendering. In \textit{Proceedings of the IEEE/CVF Conference on Computer Vision and Pattern Recognition (CVPR)}. 1440--1449.

    \item Marc Levoy and Pat Hanrahan. 1996. Light field rendering. In \textit{Proceedings of the 23rd Annual Conference on Computer Graphics and Interactive Techniques}. 31--42.

    \item Stephen Lombardi, Tomas Simon, Gabriel Schwartz, Michael Zollh{\"o}fer, Yaser Sheikh, and Jason Saragih. 2021. Mixture of volumetric primitives for efficient neural rendering. \textit{ACM Transactions on Graphics (TOG)} 40, 4 (2021), 1--13.

    \item Duane G. Merrill and Andrew S. Grimshaw. 2010. Revisiting sorting for GPGPU stream architectures. In \textit{Proceedings of the 19th International Conference on Parallel Architectures and Compilation Techniques}. 545--546.

    \item Ben Mildenhall, Pratul P. Srinivasan, Matthew Tancik, Jonathan T. Barron, Ravi Ramamoorthi, and Ren Ng. 2020. NeRF: Representing Scenes as Neural Radiance Fields for View Synthesis. In \textit{European Conference on Computer Vision (ECCV)}.

    \item Thomas M{\"u}ller, Alex Evans, Christoph Schied, and Alexander Keller. 2022. Instant Neural Graphics Primitives with a Multiresolution Hash Encoding. \textit{ACM Transactions on Graphics} 41, 4, Article 102 (July 2022), 15 pages. \url{https://doi.org/10.1145/3528223.3530127}.

    \item Eric Penner and Li Zhang. 2017. Soft 3D reconstruction for view synthesis. \textit{ACM Transactions on Graphics (TOG)} 36, 6 (2017), 1--11.

    \item Hanspeter Pfister, Matthias Zwicker, Jeroen van Baar, and Markus Gross. 2000. Surfels: Surface Elements as Rendering Primitives. In \textit{Proceedings of the 27th Annual Conference on Computer Graphics and Interactive Techniques (SIGGRAPH '00)}. ACM Press/Addison-Wesley Publishing Co., USA, 335--342. \url{https://doi.org/10.1145/344779.344936}.

    \item Christian Reiser, Songyou Peng, Yiyi Liao, and Andreas Geiger. 2021. KiloNeRF: Speeding up Neural Radiance Fields with Thousands of Tiny MLPs. In \textit{Proceedings of the IEEE/CVF International Conference on Computer Vision (ICCV)}.

    \item Liu Ren, Hanspeter Pfister, and Matthias Zwicker. 2002. Object Space EWA Surface Splatting: A Hardware Accelerated Approach to High Quality Point Rendering. \textit{Computer Graphics Forum} 21 (2002).

    \item Helge Rhodin, Nadia Robertini, Christian Richardt, Hans-Peter Seidel, and Christian Theobalt. 2015. A versatile scene model with differentiable visibility applied to generative pose estimation. In \textit{Proceedings of the IEEE International Conference on Computer Vision}. 765--773.

    \item Gernot Riegler and Vladlen Koltun. 2020. Free view synthesis. In \textit{European Conference on Computer Vision}. Springer, 623--640.

    \item Darius R{\"u}ckert, Linus Franke, and Marc Stamminger. 2022. ADOP: Approximate Differentiable One-Pixel Point Rendering. \textit{ACM Transactions on Graphics} 41, 4, Article 99 (Jul 2022), 14 pages. \url{https://doi.org/10.1145/3528223.3530122}.

    \item Miguel Sainz and Renato Pajarola. 2004. Point-based rendering techniques. \textit{Computers and Graphics} 28, 6 (2004), 869--879. \url{https://doi.org/10.1016/j.cag.2004.08.014}.

    \item Johannes Lutz Sch{\"o}nberger and Jan-Michael Frahm. 2016. Structure-from-Motion Revisited. In \textit{Proceedings of the IEEE Conference on Computer Vision and Pattern Recognition (CVPR)}.

    \item Markus Sch{\"u}tz, Bernhard Kerbl, and Michael Wimmer. 2022. Software Rasterization of 2 Billion Points in Real Time. \textit{Proceedings of the ACM on Computer Graphics and Interactive Techniques} 5, 3, Article 24 (Jul 2022), 17 pages. \url{https://doi.org/10.1145/3543863}.

    \item Vincent Sitzmann, Justus Thies, Felix Heide, Matthias Nie{\ss}ner, Gordon Wetzstein, and Michael Zollh{\"o}fer. 2019. DeepVoxels: Learning persistent 3D feature embeddings. In \textit{Proceedings of the IEEE/CVF Conference on Computer Vision and Pattern Recognition}. 2437--2446.

    \item Noah Snavely, Steven M. Seitz, and Richard Szeliski. 2006. Photo tourism: exploring photo collections in 3D. In \textit{Proceedings of SIGGRAPH}.

    \item Carsten Stoll, Nils Hasler, Juergen Gall, Hans-Peter Seidel, and Christian Theobalt. 2011. Fast articulated motion tracking using a sums of gaussians body model. In \textit{2011 International Conference on Computer Vision}. IEEE, 951--958.

    \item Cheng Sun, Min Sun, and Hwann-Tzong Chen. 2022. Direct Voxel Grid Optimization: Super-fast Convergence for Radiance Fields Reconstruction. In \textit{Proceedings of the IEEE/CVF Conference on Computer Vision and Pattern Recognition}.

    \item Towaki Takikawa, Alex Evans, Jonathan Tremblay, Thomas M{\"u}ller, Morgan McGuire, Alec Jacobson, and Sanja Fidler. 2022. Variable bitrate neural fields. In \textit{ACM SIGGRAPH 2022 Conference Proceedings}. 1--9.

    \item Towaki Takikawa, Joey Litalien, Kangxue Yin, Karsten Kreis, Charles Loop, Derek Nowrouzezahrai, Alec Jacobson, Morgan McGuire, and Sanja Fidler. 2021. Neural Geometric Level of Detail: Real-time Rendering with Implicit 3D Shapes. 2021.

    \item Ayush Tewari, Justus Thies, Ben Mildenhall, Pratul Srinivasan, Edgar Tretschk, Yifan Weng, Christoph Lassner, Vincent Sitzmann, Ricardo Martin-Brualla, Stephen Lombardi, et al. 2022. Advances in neural rendering. In \textit{Computer Graphics Forum}, Vol. 41. Wiley Online Library, 703--735.

    \item Justus Thies, Michael Zollh{\"o}fer, and Matthias Nie{\ss}ner. 2019. Deferred neural rendering: Image synthesis using neural textures. \textit{ACM Transactions on Graphics (TOG)} 38, 4 (2019), 1--12.

    \item Angtian Wang, Peng Wang, Jian Sun, Adam Kortylewski, and Alan Yuille. 2023. VoGE: A Differentiable Volume Renderer using Gaussian Ellipsoids for Analysis-by-Synthesis. In \textit{The Eleventh International Conference on Learning Representations}.

    \item Olivia Wiles, Georgia Gkioxari, Richard Szeliski, and Justin Johnson. 2020. SynSin: End-to-end view synthesis from a single image. In \textit{Proceedings of the IEEE/CVF Conference on Computer Vision and Pattern Recognition}. 7467--7477.

    \item Xiuchao Wu, Jiamin Xu, Zihan Zhu, Hujun Bao, Qixing Huang, James Tompkin, and Weiwei Xu. 2022. Scalable Neural Indoor Scene Rendering. \textit{ACM Transactions on Graphics (TOG)} (2022).

    \item Yiheng Xie, Towaki Takikawa, Shunsuke Saito, Or Litany, Shiqin Yan, Numair Khan, Federico Tombari, James Tompkin, Vincent Sitzmann, and Srinath Sridhar. 2022. Neural fields in visual computing and beyond. In \textit{Computer Graphics Forum}, Vol. 41. Wiley Online Library, 641--676.

    \item Qiangeng Xu, Zexiang Xu, Julien Philip, Sai Bi, Zhixin Shu, Kalyan Sunkavalli, and Ulrich Neumann. 2022. Point-NeRF: Point-based neural radiance fields. In \textit{Proceedings of the IEEE/CVF Conference on Computer Vision and Pattern Recognition}. 5438--5448.

    \item Yifan Wang, Felice Serena, Shihao Wu, Cengiz {\"O}ztireli, and Olga Sorkine-Hornung. 2019. Differentiable surface splatting for point-based geometry processing. \textit{ACM Transactions on Graphics (TOG)} 38, 6 (2019), 1--14.

    \item Alex Yu, Ruilong Li, Matthew Tancik, Hao Li, Ren Ng, and Angjoo Kanazawa. 2021. PlenOctrees for Real-time Rendering of Neural Radiance Fields. In \textit{Proceedings of the IEEE/CVF International Conference on Computer Vision}.

    \item Qiang Zhang, Seung-Hwan Baek, Szymon Rusinkiewicz, and Felix Heide. 2022. Differentiable Point-Based Radiance Fields for Efficient View Synthesis. In \textit{SIGGRAPH Asia 2022 Conference Papers (SA '22)}. Association for Computing Machinery, New York, NY, USA, Article 7, 12 pages. \url{https://doi.org/10.1145/3550469.3555413}.

    \item Tinghui Zhou, Shubham Tulsiani, Weilun Sun, Jitendra Malik, and Alexei A. Efros. 2016. View synthesis by appearance flow. In \textit{European Conference on Computer Vision}. Springer, 286--301.

    \item Matthias Zwicker, Hanspeter Pfister, Jeroen van Baar, and Markus Gross. 2001a. EWA volume splatting. In \textit{Proceedings Visualization, 2001. VIS '01}. IEEE, 29--538.

    \item Matthias Zwicker, Hanspeter Pfister, Jeroen van Baar, and Markus Gross. 2001b. Surface Splatting. In \textit{Proceedings of the 28th Annual Conference on Computer Graphics and Interactive Techniques (SIGGRAPH '01)}. Association for Computing Machinery, New York, NY, USA, 371--378. \url{https://doi.org/10.1145/383259.383300}.
\end{enumerate}
\endgroup


\section{外文附件}

\subsection{附录 A 梯度计算细节}

作者回顾了正文中的记号：\(\Sigma\) 与 \(\Sigma'\) 分别表示高斯在世界空间与视图空间中的协方差矩阵，\(q\) 表示旋转四元数，\(s\) 表示缩放，\(W\) 表示观察变换，而 \(J\) 表示对投影变换进行仿射近似后得到的雅可比矩阵。由此，可以先用链式法则把关于缩放和旋转的梯度写成：
\begin{equation}
    \frac{\partial \Sigma'}{\partial s}
    =
    \frac{\partial \Sigma'}{\partial \Sigma}
    \frac{\partial \Sigma}{\partial s}.
\end{equation}
\begin{equation}
    \frac{\partial \Sigma'}{\partial q}
    =
    \frac{\partial \Sigma'}{\partial \Sigma}
    \frac{\partial \Sigma}{\partial q}.
\end{equation}

将正文第 5 式中的 \(U = JW\) 代入，并把 \(\Sigma'\) 视为 \(U\Sigma U^\top\) 的左上 \(2\times2\) 对称子块，就可以得到关于 \(\Sigma\) 的偏导。随后，作者继续把三维高斯的协方差写成：
\begin{equation}
    \Sigma = RSS^\top R^\top .
\end{equation}
若令 \(M = RS\)，则又有 \(\Sigma = MM^\top\)。由于协方差矩阵及其梯度都是对称矩阵，因此可以把对 \(M\) 的求导写得更紧凑；对于尺度参数，相应的分量导数可写为
\begin{equation}
    \frac{\partial M_{ij}}{\partial s_k}
    =
    R_{ik} \quad (j = k),
    \qquad
    \frac{\partial M_{ij}}{\partial s_k}
    =
    0 \quad (j \ne k).
\end{equation}

为了求旋转的梯度，作者接着使用单位四元数 \(\mathbf{q}=(q_r, q_i, q_j, q_k)\) 到旋转矩阵的转换关系，原文式（10）如图~\ref{fig:3dgs-appendix-eq10} 所示。在此基础上，作者进一步给出了四元数四个分量对应的梯度展开式，如图~\ref{fig:3dgs-appendix-eq11} 所示；至于四元数归一化本身的梯度，作者指出它可以直接推得。

\begin{figure}[H]
    \centering
    \includegraphics[width=0.96\textwidth]{translation/3dgs-appendix-eq10.png}
    \caption{附录式（10）：单位四元数到旋转矩阵的转换关系。}
    \label{fig:3dgs-appendix-eq10}
\end{figure}

\begin{figure}[H]
    \centering
    \includegraphics[width=0.96\textwidth]{translation/3dgs-appendix-eq11.png}
    \caption{附录式（11）：四元数各个分量对应的梯度展开式。}
    \label{fig:3dgs-appendix-eq11}
\end{figure}

\subsection{附录 B 优化与致密化算法}

附录 B 将正文中的训练流程整理成算法形式。整体上，方法从 SfM 点集出发初始化高斯的位置，再初始化协方差、颜色和透明度，随后不断采样训练视图、执行光栅化、计算损失并进行 Adam 更新；在预定的细化迭代上，再执行裁剪、克隆和分裂等致密化操作。该流程的中文化整理如下：

\begin{algorithm}[H]
    \caption{优化与高斯致密化算法}
    \small
    \begin{algorithmic}[1]
        \STATE 输入：训练图像宽度 \(w\)、高度 \(h\)
        \STATE \(M \leftarrow \) SfM 点集
        \STATE \((S, C, A) \leftarrow \texttt{InitAttributes}()\)
        \STATE \(i \leftarrow 0\)
        \WHILE{算法尚未收敛}
            \STATE \((V, I) \leftarrow \texttt{SampleTrainingView}()\)
            \STATE \(\hat{I} \leftarrow \texttt{Rasterize}(M, S, C, A, V)\)
            \STATE \(\mathcal{L} \leftarrow \texttt{Loss}(I, \hat{I})\)
            \STATE \((M, S, C, A) \leftarrow \texttt{Adam}(\nabla \mathcal{L})\)
            \IF{\(\texttt{IsRefinementIteration}(i)\)}
                \FORALL{高斯 \((\mathbf{x}, \Sigma, \mathbf{c}, \alpha)\)}
                    \IF{\(\alpha < \epsilon_{\alpha}\) 或 \(\texttt{IsTooLarge}(\Sigma, V)\)}
                        \STATE \texttt{RemoveGaussian}()
                    \ENDIF
                    \IF{\(\| \nabla_{\mathbf{p}} \mathcal{L} \| > \tau_{\mathbf{p}}\)}
                        \IF{\(\| \Sigma \| > r_g\)}
                            \STATE \texttt{SplitGaussian}\((\mathbf{x}, \Sigma, \mathbf{c}, \alpha)\)
                        \ELSE
                            \STATE \texttt{CloneGaussian}\((\mathbf{x}, \Sigma, \mathbf{c}, \alpha)\)
                        \ENDIF
                    \ENDIF
                \ENDFOR
            \ENDIF
            \STATE \(i \leftarrow i + 1\)
        \ENDWHILE
    \end{algorithmic}
\end{algorithm}

\subsection{附录 C 光栅化器细节}

附录 C 首先补充了排序步骤的实现细节。作者的设计假设场景中存在大量小尺度泼溅，因此选择在每一帧开始时，用一次基数排序完成全部泼溅的排序。具体做法是：先把屏幕划分为 \(16\times16\) 的小块，再把每个泼溅复制到它所覆盖的每一个小块中，从而形成小块与泼溅实例的对应关系。虽然这会适度增加待处理的高斯数量，但由于控制流更加简单，并且可以充分利用 GPU 基数排序的并行性，整体仍然是划算的。每个泼溅实例都会获得一个最高可达 \(64\) 位的键值，其中低 \(32\) 位编码投影深度，高位则编码其所在小块编号。这样一来，就能通过一次全局排序并行地完成深度次序与小块归属的组织。

完成排序之后，作者再根据相同小块编号在有序数组中的起止位置，为每个小块生成对应的高斯列表。随后，前向渲染阶段只需要在每个小块内部按照既定顺序执行混合即可。附录中的 CPU 软件光栅化流程可整理为下面的算法：

\begin{algorithm}[H]
    \caption{三维高斯的 CPU 软件光栅化流程}
    \small
    \begin{algorithmic}[1]
        \STATE 输入：待光栅化图像的宽度 \(w\)、高度 \(h\)
        \STATE 输入：世界空间中的高斯均值与协方差 \((M, S)\)
        \STATE 输入：高斯颜色 \(C\)、透明度 \(A\) 以及当前相机视图配置 \(V\)
        \STATE \texttt{CullGaussian}\((\mathbf{p}, V)\)，执行视锥裁剪
        \STATE \((M', S') \leftarrow \texttt{ScreenSpaceGaussians}(M, S, V)\)
        \STATE \(T \leftarrow \texttt{CreateTiles}(w, h)\)
        \STATE \((L, K) \leftarrow \texttt{DuplicateWithKeys}(M', T)\)
        \STATE \texttt{SortByKeys}\((K, L)\)
        \STATE \(R \leftarrow \texttt{IdentifyTileRanges}(T, K)\)
        \STATE \(I \leftarrow 0\)，初始化画布
        \FORALL{小块 \(t \in T\)}
            \FORALL{小块 \(t\) 中的像素 \(i\)}
                \STATE \(r \leftarrow \texttt{GetTileRange}(R, t)\)
                \STATE \(I[i] \leftarrow \texttt{BlendInOrder}(i, L, r, K, M', S', C, A)\)
            \ENDFOR
        \ENDFOR
        \STATE \textbf{return} \(I\)
    \end{algorithmic}
\end{algorithm}

附录 C 还进一步讨论了数值稳定性。在反向传播中，作者需要通过不断用各高斯的 \(\alpha\) 去除前向传播中累积的不透明度，来恢复梯度计算所需的中间不透明度值。如果直接这样做，会出现除以零等数值不稳定问题。为此，作者在前向与反向两个阶段都跳过 \(\alpha < \epsilon\) 的混合更新，并将 \(\alpha\) 的上界截断到 \(0.99\)。此外，在某个高斯被纳入前向光栅化之前，还会预先计算若把它加入当前累积结果后得到的不透明度，并在该值将超过 \(0.9999\) 之前停止前向的从前到后混合。

\subsection{附录 D 各场景误差指标}

附录 D 汇总了论文评测中使用的逐场景误差指标。表 4 到表 9 给出了所有比较方法在真实场景上的 SSIM、PSNR 与 LPIPS 数值；其中既包括直接摘自 Mip-NeRF360 原论文的数值，也包括作者为了生成本文可视化结果而重新运行得到的数值。作者最后还给出了其在完整 Mip-NeRF360 数据集上的平均值：PSNR 为 \(27.58\)，SSIM 为 \(0.790\)，LPIPS 为 \(0.240\)。

\begin{figure}[H]
    \centering
    \includegraphics[width=0.98\textwidth]{translation/3dgs-appendix-tables-1.png}
    \caption{附录表 4 至表 6：Mip-NeRF360 场景上的 SSIM、PSNR 与 LPIPS 指标。带匕首符号的结果直接摘自原论文。}
\end{figure}

\begin{figure}[H]
    \centering
    \includegraphics[width=0.92\textwidth]{translation/3dgs-appendix-tables-2.png}
    \caption{附录表 7 至表 9：Tanks\&Temples 与 Deep Blending 场景上的 SSIM、PSNR 与 LPIPS 指标。}
\end{figure}

\begingroup
    \linespreadsingle{}
    \printbibliography[title={本章引用文献}]
\endgroup
